\documentclass[12pt,a4paper]{article}
\usepackage[utf8]{inputenc}
\usepackage[spanish,es-noquoting]{babel}
\usepackage[T1]{fontenc}
\usepackage{lmodern}
\usepackage{geometry}
\geometry{top=2cm, bottom=2cm, left=2.5cm, right=2.5cm}
\usepackage{xcolor}
\usepackage{listings}
\usepackage{enumitem}
\usepackage{fancyhdr}
\usepackage{hyperref}
\usepackage{etoolbox}
\usepackage{needspace}
\usepackage{tikz}
\usetikzlibrary{arrows.meta, positioning, calc}

\definecolor{codegreen}{rgb}{0,0.6,0}
\definecolor{backcolour}{rgb}{0.95,0.95,0.92}

\lstset{
    language=C,
    backgroundcolor=\color{backcolour},
    commentstyle=\color{codegreen},
    keywordstyle=\color{blue},
    basicstyle=\ttfamily\small,
    breaklines=true,
    frame=single,
    numbers=left,
    tabsize=4,
    extendedchars=true,
    showstringspaces=false,
    inputencoding=utf8,
    literate={á}{{\'a}}1 {é}{{\'e}}1 {í}{{\'i}}1 {ó}{{\'o}}1 {ú}{{\'u}}1 {ñ}{{\~n}}1 {¡}{{\textexclamdown}}1 {¿}{{\textquestiondown}}1
}

\usepackage{graphicx}
\graphicspath{{../../../assets/}{../../assets/}{assets/}}
\BeforeBeginEnvironment{lstlisting}{\Needspace{4cm}}

\newcommand{\unilogo}{%
  \IfFileExists{../../../assets/logo_unimagdalena.png}%
    {\includegraphics[height=1.6cm]{logo_unimagdalena.png}}%
    {\IfFileExists{../../assets/logo_unimagdalena.png}%
      {\includegraphics[height=1.6cm]{logo_unimagdalena.png}}%
      {\IfFileExists{assets/logo_unimagdalena.png}%
        {\includegraphics[height=1.6cm]{logo_unimagdalena.png}}%
        {}}}%
}

\setlength{\headheight}{15pt}
\pagestyle{fancy}
\fancyhf{}
\lhead{Monitorias Sistemas Operativos 2026-II}
\rhead{Capítulo 2: Colas de Mensajes}
\cfoot{\thepage}
\fancypagestyle{plain}{%
    \fancyhf{}%
    \lhead{Monitorias Sistemas Operativos 2026-II}%
    \rhead{Capítulo 2: Colas de Mensajes}%
    \cfoot{\thepage}%
}
\fancypagestyle{titlepage}{%
    \fancyhf{}%
    \cfoot{\thepage}%
    \renewcommand{\headrulewidth}{0pt}%
}

\title{
\begin{tikzpicture}[remember picture, overlay]
    \node[anchor=west, inner sep=0pt, xshift=2.5cm, yshift=-2.0cm] at (current page.north west) {
        \unilogo
    };
    \node[anchor=east, inner sep=0pt, xshift=-2.5cm, yshift=-2.0cm] at (current page.north east) {
        \large\textbf{2026-II}
    };
\end{tikzpicture}%
\textbf{Taller Propuesto: Análisis Estadístico con Colas de Mensajes}}
\author{Malak Sanchez \\ \small Monitorias Sistemas Operativos}
\date{\today}

\begin{document}

\maketitle
\thispagestyle{titlepage}

\section*{Introducción}
Este taller plantea un problema de procesamiento estadístico concurrente que debe resolverse con colas de mensajes System V. El estudiante deberá leer y validar datos almacenados en un archivo, diseñar un protocolo entre un proceso coordinador y varios trabajadores, combinar resultados parciales e implementar la solución en C.

\section*{Problema Propuesto}

\subsection*{Contexto}
Un laboratorio recibe archivos con mediciones enteras provenientes de distintos sensores. Cada archivo puede contener miles de valores, por lo que se desea repartir el análisis entre varios procesos. Un proceso coordinador debe leer las mediciones, crear un grupo de trabajadores y distribuir los datos. Cada trabajador calcula un resumen estadístico parcial y lo devuelve al coordinador.

La comunicación debe realizarse mediante \textbf{una sola cola de mensajes System V}. No se permite utilizar tuberías, memoria compartida, archivos temporales ni espera activa para intercambiar tareas o resultados. Solo el proceso padre puede leer el archivo de entrada.

\subsection*{Enunciado}
Desarrolle un programa llamado \texttt{analizador\_estadistico} que reciba por línea de comandos la cantidad de trabajadores y la ruta de un archivo de mediciones. El proceso padre debe leer y validar el archivo antes de crear los hijos. Luego debe distribuir los valores entre los trabajadores en orden circular mediante una cola de mensajes.

Cada trabajador debe mantener acumuladores locales y, al terminar sus tareas, enviar un único resumen parcial al padre. El padre debe combinar los resúmenes, calcular las estadísticas globales, esperar la terminación de todos los hijos y eliminar la cola.

\subsection*{Ejecución y formato del archivo}
El programa se ejecutará con la siguiente forma:

\begin{lstlisting}[language={}, numbers=none]
./analizador_estadistico <trabajadores> <archivo>
\end{lstlisting}

La primera línea del archivo contiene la cantidad $N$ de mediciones. A continuación aparecen exactamente $N$ enteros separados por espacios o saltos de línea.

\begin{lstlisting}[language={}, numbers=none]
8
12 7 -3 18 5 21 10 4
\end{lstlisting}

Considere las siguientes restricciones:
\begin{itemize}
    \item $1 \leq N \leq 10000$.
    \item Cada medición debe pertenecer al intervalo $[-100000,100000]$.
    \item La cantidad de trabajadores debe estar entre 1 y $\min(16,N)$.
    \item El programa debe rechazar archivos inexistentes, valores no enteros, cantidades fuera de rango, datos faltantes y datos adicionales.
    \item Una entrada inválida debe producir un mensaje de error y una terminación controlada sin dejar recursos IPC activos.
\end{itemize}

\subsection*{Análisis estadístico requerido}
Cada trabajador debe calcular sobre sus valores asignados: cantidad de elementos, suma, suma de cuadrados, mínimo, máximo y cantidad de números primos. El padre combina los resúmenes parciales y calcula:
\begin{itemize}
    \item cantidad total de mediciones;
    \item mínimo y máximo globales;
    \item media aritmética $\overline{x} = \frac{\sum x_i}{N}$;
    \item varianza poblacional $\sigma^2 = \frac{\sum x_i^2}{N} - \overline{x}^{,2}$;
    \item desviación estándar poblacional $\sigma = \sqrt{\sigma^2}$;
    \item cantidad y porcentaje de números primos.
\end{itemize}

El programa debe imprimir las cantidades enteras sin decimales y las medidas reales con dos cifras decimales. Los valores menores que 2 no se consideran primos.

\subsection*{Requisitos funcionales}

\begin{enumerate}[itemsep=0.35em]
    \item Leer y validar completamente el archivo en el proceso padre antes de crear la cola y los trabajadores.
    \item Crear una cola privada con \texttt{msgget(IPC\_PRIVATE, IPC\_CREAT | 0600)} antes de invocar \texttt{fork()}.
    \item Crear $P$ trabajadores. El trabajador $i$ recibe únicamente mensajes de tipo $i$, con tipos numerados desde 1.
    \item Distribuir los trabajos en orden circular: el trabajo $j$ se envía al trabajador $(j \bmod P) + 1$.
    \item Enviar cada tarea como una estructura con identificador de trabajo y valor entero.
    \item Enviar un mensaje de terminación a cada trabajador después de distribuir todos los valores.
    \item Hacer que cada trabajador envíe exactamente un resumen parcial con un tipo común reservado. El cuerpo debe contener su PID y todos los acumuladores necesarios.
    \item Recibir $P$ resúmenes, combinarlos y mostrar las estadísticas globales solicitadas.
    \item Esperar a todos los hijos con \texttt{waitpid()} y eliminar la cola con \texttt{msgctl(IPC\_RMID)}.
    \item Validar argumentos y comprobar los valores de retorno de todas las llamadas al sistema.
\end{enumerate}

\subsection*{Restricciones técnicas}
\begin{itemize}[itemsep=0.25em]
    \item Utilice una única cola de mensajes para tareas, terminación y resúmenes parciales.
    \item El primer campo de toda estructura enviada debe ser un \texttt{long} positivo que represente el tipo de mensaje.
    \item El tamaño enviado a \texttt{msgsnd()} y \texttt{msgrcv()} debe corresponder solo al cuerpo del mensaje.
    \item Los trabajadores deben esperar con \texttt{msgrcv()}; no se permite sondear la cola con ciclos de espera activa.
    \item Use tipos capaces de almacenar de forma segura la suma y la suma de cuadrados dentro de los límites establecidos.
    \item El proceso padre debe ser el único responsable de eliminar la cola.
    \item La solución debe compilar sin advertencias con \texttt{-Wall -Wextra -Wpedantic} y enlazar la biblioteca matemática con \texttt{-lm}.
\end{itemize}

\subsection*{Arquitectura requerida}

\begin{center}
\begin{tikzpicture}[
    >=Stealth,
    node distance=1.4cm,
    proc/.style={draw, thick, rounded corners, fill=blue!8, minimum width=2.5cm, minimum height=0.9cm, align=center},
    queue/.style={draw, thick, fill=yellow!20, minimum width=3.2cm, minimum height=1.1cm, align=center}
]
    \node[proc] (parent) {Padre\\coordinador};
    \node[draw, fill=gray!10, above=0.8cm of parent, align=center] (file) {Archivo\\$N$ enteros};
    \node[queue, right=2.2cm of parent] (queue) {Cola System V\\tipos 1, 2, $\ldots$, $P$ y 1000};
    \node[proc, above right=0.2cm and 2.2cm of queue] (w1) {Trabajador 1\\\texttt{msgrcv(..., 1, 0)}};
    \node[proc, right=2.2cm of queue] (w2) {Trabajador 2\\\texttt{msgrcv(..., 2, 0)}};
    \node[proc, below right=0.2cm and 2.2cm of queue] (wn) {Trabajador $P$\\\texttt{msgrcv(..., P, 0)}};

    \draw[->, thick] (file) -- node[right] {lee} (parent);
    \draw[->, thick] ([yshift=5pt]parent.east) -- node[above] {tareas} ([yshift=5pt]queue.west);
    \draw[<-, thick, violet!70!black] ([yshift=-5pt]parent.east) -- node[below] {resúmenes} ([yshift=-5pt]queue.west);
    \draw[<->, thick] (queue.north east) -- (w1.west);
    \draw[<->, thick] (queue.east) -- (w2.west);
    \draw[<->, thick] (queue.south east) -- (wn.west);
\end{tikzpicture}
\end{center}

\newpage
\subsection*{Preguntas de análisis}
\thispagestyle{fancy}
\begin{enumerate}
    \item ¿Por qué la distribución circular utiliza $P$ y no $N$ en la operación módulo? Justifique su respuesta.
    \item ¿Por qué es preferible que solo el padre lea el archivo en lugar de que cada trabajador vuelva a abrirlo?
    \item Explique cómo se combinan los mínimos, máximos, sumas y sumas de cuadrados parciales sin conservar nuevamente todos los datos.
    \item ¿Por qué la suma y la suma de cuadrados no deberían almacenarse en variables de tipo \texttt{int}?
    \item ¿Cómo evita el protocolo que un trabajador consuma un resumen destinado al padre?
    \item ¿Qué podría ocurrir si el padre ejecutara \texttt{waitpid()} antes de recibir los resúmenes y la cola llegara a llenarse?
    \item ¿Cómo garantiza su implementación que solamente el proceso propietario elimine la cola?
    \item Calcule cuántos mensajes se intercambian en función de $N$ y $P$. Compare este protocolo con uno donde cada trabajador responda por cada valor procesado.
    \item ¿Cómo podría detectar el padre que un trabajador terminó inesperadamente antes de enviar su resumen, sin quedar bloqueado indefinidamente?
    \item La fórmula de la varianza puede producir un valor negativo muy cercano a cero debido al redondeo. ¿Cómo debería tratarse ese caso?
\end{enumerate}

\end{document}
